% Source: https://artofproblemsolving.com/community/c3h2067143_coordinate_plane_problem
% Stable identity: aops-topic-2067143-q01
For any line through $P$, the distance from $O$ to the line is at most $OP$: the perpendicular from $O$ to the line is no longer than the segment $OP$ to a point of the line. Equality holds exactly when $OP$ is itself the perpendicular, i.e.\ when the line is perpendicular to $OP$ at $P$.

\begin{asy}
usepackage("amsmath");
size(6cm);
defaultpen(fontsize(10pt));
pair O=(0,0), P=(3,5);
draw((-1,0)--(13,0),gray(0.6),Arrow); draw((0,-1)--(0,8),gray(0.6),Arrow);
draw((-0.5,(34+1.5)/5)--(12,(34-36)/5),blue+linewidth(0.9));
draw(O--P,red+linewidth(0.9));
pair u=unit(P-O), v=unit((5,-3));
draw(P-0.4*u--P-0.4*u+0.4*v--P+0.4*v);
pair Q=(8,2); draw(O--Q,dashed); pair Hq=Q; dot(Q);
dot(O);dot(P);
label("$O$",O,SW); label("$P(3,5)$",P,N); label("$\ell:\ 3x+5y=34$",(9.5,(34-28.5)/5),NE,blue); label("$Q$",Q,NE);
label("\mbox{}",(min(currentpicture,user=true)+max(currentpicture,user=true))/2,black); // keeps the text after the figure black when LaTeX typesets the labels inline
\end{asy}

The slope of $OP$ is $\dfrac{5}{3}$, so $\ell$ has slope $-\dfrac{3}{5}$ and passes through $(3,5)$:
\[
y-5=-\dfrac{3}{5}(x-3)\iff \boxed{3x+5y=34}.
\]
Its distance from $O$ is $\dfrac{34}{\sqrt{34}}=\sqrt{34}=OP$, as expected. (For any other point $Q$ of $\ell$, $OQ>OP$, as in the dashed segment.)
